% Draft paper. Written in English (more publishable than French for a transport
% / environmental-emissions venue); the validation protocol it references is in
% French in validation-experimentale.md.
\documentclass[11pt,a4paper]{article}
\usepackage[utf8]{inputenc}
\usepackage[T1]{fontenc}
\usepackage{lmodern}
\usepackage[margin=2.5cm]{geometry}
\usepackage{graphicx}
\usepackage{booktabs}
\usepackage{amsmath}
\usepackage{amssymb}
\usepackage{siunitx}
\usepackage{hyperref}
\usepackage{authblk}
\usepackage{caption}

\title{Cruise speed versus vehicle class on a short, steep mountain road:\\
a transparent fuel and particulate-emission simulator\\
with paired field validation}

\author[1]{M.~Dutour}
\affil[1]{Independent --- \texttt{simulateur-vitesse-pollution-montagne}, Plateau d'Assy, France}

\date{\today}

\begin{document}
\maketitle

\begin{abstract}
We describe a fully transparent, browser-side simulator that estimates
fuel consumption, \ce{CO2}, exhaust and non-exhaust particulate matter
for an internal-combustion-engine vehicle traversing the Super~U~Passy
$\leftrightarrow$ Maison m\'edicale du Plateau d'Assy mountain road
(\SI{10.235}{km} one-way, \SI{460}{m} elevation change). The model
combines an OSM-derived geometry, a kinematic speed profile bounded by
posted limits, curve geometry and engine power, a Willans-line fuel
model with idle and driveline losses, EMEP/EEA non-exhaust factors
scaled by mass, and a kinetic-energy-based brake-PM model after the
engine-braking share is removed. Every coefficient is sourced and
exposed in the user interface. We propose a paired within-subject
field campaign that targets the \emph{difference} between two cruise
speeds rather than absolute values, designed so the simulator can be
falsified or accepted as a comparison tool even in the presence of
moderate absolute bias. Pre-registered effect sizes and acceptance
intervals are given. Implications for advisory speed signage on short,
steep roads are discussed.
\end{abstract}

\section{Introduction}

Air-quality policy on small mountain roads typically inherits speed
limits and emission assumptions from urban or motorway contexts. Yet
short, steep links exhibit two effects that make those assumptions
unreliable: the fraction of total trip energy spent climbing is large,
so any aero-versus-gravity speed-sensitivity intuition flips; and a
non-trivial share of the kilometer-averaged particulate emission
comes from non-exhaust sources, whose dependence on speed is itself
contested~\cite{beddows2021,emep2023ldv}.

We present a simulator built specifically for a single
\SI{10.235}{km} road on the Plateau d'Assy massif in the Haute-Savoie
prealps. The choice is deliberate: the road carries daily medical
traffic, has a posted limit that varies between \SI{30}{km/h} and
\SI{90}{km/h}, climbs \SI{460}{m} over its length, and includes
nineteen curves with radii as small as \SI{25}{m}. Public discussion
about an advisory speed limit on this link motivated the work.

The simulator is deliberately scoped so that:

\begin{enumerate}
  \item every coefficient is traceable to a public source visible
        in the user interface;
  \item the code is short (\textasciitilde\num{1700} lines, no
        dependencies, runs entirely in the browser);
  \item the user controls all relevant degrees of freedom (cruise
        speed, vehicle profile, comfort acceleration, cold-start
        budget, direction, round trip);
  \item the output that matters --- the \emph{difference} between
        two scenarios --- is presented as the primary metric.
\end{enumerate}

This paper documents the model, sketches a paired field-validation
campaign aimed at the comparison axis, and lays out the
pre-registered statistical machinery for accepting or falsifying it.

\section{Route and study site}

The route follows D~39, D~43 and D~13 between the Super~U Passy
gas station (\ang{45;55;9.2} N, \ang{6;42;16.1} E, \SI{578}{m}) and
the Maison m\'edicale du Plateau d'Assy (\ang{45;56;22.1} N,
\ang{6;42;26.8} E, \SI{1039}{m}). Geometry is taken from
OpenStreetMap; elevations from the European EU-DEM at \SI{25}{m}
spacing; \texttt{maxspeed} tags are used where present, with
fallbacks of \SI{50}{km/h} (urban) and \SI{90}{km/h} (rural)
otherwise. The route includes nineteen labelled curves whose radii
were estimated from the OSM polyline.

\section{Model}

The simulator builds a position-discretised speed profile
$v(s)$ on $s \in [0, L]$ with $L$ the route length (one-way or doubled
for round trip), then integrates a longitudinal force balance
\cite{liu2017} along $s$ to obtain energy and emission totals.

\subsection{Speed profile}

For each station~$s$ the local speed cap is

\begin{equation}
v_{\text{cap}}(s) = \min\!\left( v_{\text{cruise}},\, v_{\text{lim}}(s),\,
  \min_{c \in \mathcal{C}} v_{\text{curve}}(s, c) \right),
\end{equation}

where $v_{\text{lim}}$ is the posted limit and $v_{\text{curve}}(s,c)
= \sqrt{a_{\text{lat}} R_c}$ for curve $c$ of radius $R_c$. The curve
cap holds across the full arc length $R_c \, \theta_c$, not only at
the apex. A parabolic approach lookahead with
$a_{\text{decel}} = a_{\text{comfort}} + g \sin\theta(s)$ ensures the
deceleration is consistent with the slope.

A forward pass clamps acceleration by both the comfort target and the
engine power available at the local slope and speed. We use the
canonical constant-torque-below-base-speed / constant-power-above
form~\cite{heywood2018}:

\begin{equation}
F_{\text{avail}}(v, h) = \frac{P_{\max} \,\eta_{\text{dt}} \, \rho(h)/\rho_0}
                              {\max(v_{P_{\max}}, v)}.
\end{equation}

Air density follows the ISA troposphere
$\rho(h) = \rho_0 (1 - L_h h)^{4.2559}$ with $L_h =
\SI{2.2557e-5}{m^{-1}}$. The trip starts and ends at rest, and a round
trip pauses at the turnaround.

\subsection{Longitudinal dynamics}

For each segment $[s_a, s_b]$ we compute

\begin{align}
F_{\text{roll}} &= C_{\text{rr}}(\langle v^2 \rangle) \, m g \cos\theta, \\
F_{\text{aero}} &= \tfrac{1}{2} \rho(h) C_d A \,\langle v^2 \rangle, \\
F_{\text{grade}} &= m g \sin\theta, \\
F_{\text{inertia}} &= m \, \frac{v_b^2 - v_a^2}{2(s_b - s_a)},
\end{align}

with $\langle v^2 \rangle = (v_a^2 + v_b^2)/2$ being the spatial mean
of $v^2$ for constant-acceleration kinematics --- the trapezoidal
approximation $\bar{v}^2$ is biased by Jensen's inequality on launch
and stop segments. Rolling resistance follows
$C_{\text{rr}}(v) = C_{\text{rr},0} (1 + (v / v_{\text{rr}})^2)$ with
$v_{\text{rr}} = \SI{100}{m/s}$, calibrated to the ISO~28580 ratio
$C_{\text{rr}}(\SI{130}{km/h}) / C_{\text{rr}}(\SI{50}{km/h}) \approx
1.10$.

\subsection{Fuel and \texorpdfstring{\ce{CO2}}{CO2}}

Fuel follows a Willans line~\cite{guzzella2013}

\begin{equation}
\dot m_{\text{fuel}} =
\frac{P_{\text{wheel}}}{\eta_{\text{dt}} \, \eta_{\text{ind}} \, h_u}
+ k_{\text{idle}} \, V_d \, \mathbb{1}[\text{not in DFCO}],
\end{equation}

with $\eta_{\text{ind}} = 0.40$, $\eta_{\text{dt}} = 0.85$, $k_{\text{idle}}
= \SI{0.21}{g.s^{-1}.L^{-1}}$ of displacement, and DFCO disengaged
below \SI{25}{km/h}. \ce{CO2} uses the ADEME~Base~Carbone E10
combustion factor of \SI{2.21}{kg/L}~\cite{ademe2023}.

A cold-start window of $T_{\text{cs}}$ seconds (default \SI{60}{s},
user-controlled) applies a $+30\,\%$ fuel surcharge and a
$\times 7$ catalyst-light-off PM penalty to the fraction of fuel
burned during the first $T_{\text{cs}}$ seconds (doubled for round
trips since the engine cools at the turnaround), per Weilenmann
\textit{et al.}~\cite{weilenmann2009}.

\subsection{Tailpipe PM}

Tailpipe PM follows the EMEP/EEA Tier~3 fuel-mass approach: total PM
is the fuel mass burned multiplied by a per-vehicle factor that
depends on the emissions tier (Note Euro~5: \SI{25}{mg/kg}; Ram
Tier~2 Bin~5: \SI{90}{mg/kg}; X5~Euro~3: \SI{120}{mg/kg})
\cite{emep2023si}.

\subsection{Non-exhaust PM}

Tyre and road wear use EMEP/EEA TSP factors of \SI{0.0107}{g/km} and
\SI{0.015}{g/km} respectively, scaled linearly by vehicle mass per
Beddows and Harrison~\cite{beddows2021}. Brake wear uses an
energy-based model: pad work is multiplied by \SI{0.0102}{g.MJ^{-1}}
of TSP, calibrated against the Hagino dynamometer
data~\cite{hagino2016}, and the engine-braking share
$F_{\text{eb}} = k_{\text{eb}} V_d v$ is removed before charging the
residual to the friction brakes. PM10 and PM2.5 fractions follow
EMEP~\cite{emep2023nonexhaust}.

\subsection{Diagnostics}

Two quantities not present in standard EMEP rollups are surfaced:

\begin{itemize}
  \item the \emph{power-limited distance} where
        $a_{\text{powered}}(v, \theta, \rho) < 0$, indicating
        segments where the engine cannot maintain speed at the
        local slope --- a check on the heavy-vehicle, steep-grade
        regime;
  \item per-segment brake-wear hotspots, derived from the same
        per-MJ factor as the trip total, used to colour the speed
        profile chart.
\end{itemize}

\section{Validation experiment}

The model is intended for relative comparisons between cruise speeds.
We pre-register hypotheses on the \emph{difference} between
\SI{50}{km/h} and \SI{90}{km/h} round trips (Sec.~\ref{sec:hypotheses})
and design a paired within-subject campaign whose statistical power
is dominated by between-trip noise rather than absolute model bias.
Full protocol in
\href{https://github.com/mathieudutour/simulateur-vitesse-pollution-montagne/blob/main/validation-experimentale.md}{\texttt{validation-experimentale.md}}.

\subsection{Pre-registered hypotheses}\label{sec:hypotheses}

\begin{description}
  \item[H1.] $\Delta_{\text{fuel}} = X_{90} - X_{50}$ on a round trip
        is positive, in $[+0.15, +0.40]$~L for the Note,
        $[+0.30, +0.80]$~L for the heavier profiles.
  \item[H2.] $\Delta_{\text{time}}$ is negative, in $[-9, -5]$~min,
        vehicle-independent.
  \item[H3.] $\Delta_{E_{\text{brake}}}$ is positive; the ratio
        $\Delta_{E_{\text{brake}}}^{\text{SUV}} /
         \Delta_{E_{\text{brake}}}^{\text{Note}} > 1.5$.
  \item[H4.] Power-limited distance is strictly positive on the
        climb for the Note (\SI{58}{W/kg}) and zero for the X5
        (\SI{118}{W/kg}) and Ram (\SI{107}{W/kg}); the small car,
        not the heavy ones, runs out of headroom on the steepest
        segments.
\end{description}

\subsection{Design and analysis}

We run $N=6$ paired round trips per (vehicle, cruise speed) cell
across two days separated by a week, in a Latin-square-randomised
order. Fuel is measured by tank-full to tank-full at the same pump.
Time and instantaneous speed come from a \SI{1}{Hz} GPS log. Brake
energy is proxied by IR brake-disc temperature change. PM is measured
only if a portable PEMS is available; the rest of the campaign holds
without it.

For every cell we form pairs $(\Delta_{\text{meas},i},
\Delta_{\text{sim},i})$ and run (a) a sign test on
$\Delta_{\text{meas}} - \Delta_{\text{sim}}$, (b) a paired
log-ratio test, and (c) a paired regression
$\Delta_{\text{meas}} = \alpha + \beta \Delta_{\text{sim}} +
\varepsilon$. Acceptance criterion: $\beta \in [0.7, 1.4]$, $\alpha$
indistinguishable from~$0$ at \SI{90}{\percent}, sign-test consistent.
A 90~\% confidence interval is computed by paired bootstrap rather
than Student.

\section{Results}\label{sec:results}

Empirical results from the campaign will be inserted here. The
expected qualitative shape is shown schematically in
Fig.~\ref{fig:expected}. Three observations should hold if the model
is to be useful as a public-facing tool:

\begin{enumerate}
  \item For the Note, fuel difference between \SI{50}{km/h} and
        \SI{90}{km/h} on the round trip is dominated by aero on the
        descent (where the slider acts as a soft cap) rather than
        traction on the climb (where engine power binds). The
        simulator should reproduce this asymmetry --- climb fuel is
        relatively insensitive to cruise choice.
  \item Brake energy and brake-PM differences scale super-linearly
        with mass for the same cruise-speed delta: more KE to
        dissipate when entering each curve.
  \item The cold-start budget reshapes the cross-vehicle PM ranking
        on the round trip. Without it, the EMEP fuel-mass factor
        dominates; with it, an Euro~3 SUV that cold-starts twice
        moves clearly ahead of an Euro~5 small car.
\end{enumerate}

\begin{figure}[h]
  \centering
  \includegraphics[width=0.7\linewidth]{example-image-a}
  \caption{Schematic: predicted $\Delta$(fuel, time, brake~PM) between
  \SI{50}{km/h} and \SI{90}{km/h} round trips, by vehicle profile.
  Replace with measured data after the field campaign.}
  \label{fig:expected}
\end{figure}

\section{Discussion}

The paired design intentionally trades absolute accuracy for
comparison robustness. A simulator that is, say, \SI{15}{\percent}
optimistic on absolute fuel can still be a credible comparison tool
provided its \emph{deltas} match measurement to within sampling
noise. The validation protocol is built around this trade-off.

Three limits deserve emphasis. First, the model has a single fixed
gear and uses a linear-in-vehicle-speed engine-braking term; downhill
descents in low gear are not represented and may show systematically
high friction-brake PM. Second, particle number (PN) is not modelled,
so direct-injection gasoline regulatory compliance cannot be inferred
from the output. Third, the cold-start budget is uniform in time,
whereas the real catalyst light-off curve is exponential; the total
penalty is correctly captured but the spatial distribution along the
route is not.

\section{Conclusion}

The simulator is, by construction, easier to falsify on relative
than on absolute terms; the protocol in
\href{https://github.com/mathieudutour/simulateur-vitesse-pollution-montagne/blob/main/validation-experimentale.md}{\texttt{validation-experimentale.md}}
is sized to do exactly that for the three vehicle profiles already in
the tool. If the paired comparisons hold, the simulator can be used to
inform advisory speed discussions on the Plateau d'Assy road, with
honest acknowledgement of the absolute-value uncertainty. If they do
not, the discrepancies will localise --- per-vehicle, per-direction,
per-quantity --- which coefficient or sub-model needs revision.

\paragraph{Reproducibility.} Code, route data and constants are
available at\\
\url{https://github.com/mathieudutour/simulateur-vitesse-pollution-montagne}.
The branch \texttt{claude/review-model-flaws-f6U7f} contains the model
described in this paper.

\paragraph{Acknowledgements.} OpenStreetMap and the EEA EU-DEM project
provide the route geometry and elevations.

\bibliographystyle{plain}
\begin{thebibliography}{99}

\bibitem{liu2017}
F.~Liu and Z.~Li, \emph{Modeling longitudinal vehicle dynamics on
graded roads}, Energies, vol.~10, no.~5, p.~700, 2017.

\bibitem{heywood2018}
J.~B. Heywood, \emph{Internal Combustion Engine Fundamentals},
2nd~ed., McGraw-Hill, 2018.

\bibitem{guzzella2013}
L.~Guzzella and A.~Sciarretta, \emph{Vehicle Propulsion Systems:
Introduction to Modeling and Optimization}, 3rd~ed., Springer, 2013.

\bibitem{ademe2023}
ADEME, \emph{Base Carbone, facteurs d'\'emission combustion SP95-E10},
2023, \url{https://base-carbone.ademe.fr/}.

\bibitem{weilenmann2009}
M.~Weilenmann, J.-Y. Favez, R.~Alvarez,
\emph{Cold-start emissions of modern passenger cars at different
low-ambient temperatures and their evolution over vehicle legislation
categories}, Atmospheric Environment, vol.~43, no.~15, pp.~2419--2429,
2009.

\bibitem{emep2023ldv}
European Environment Agency, \emph{EMEP/EEA Air Pollutant Emission
Inventory Guidebook 2023, Part B 1.A.3.b: Road transport},
\url{https://www.eea.europa.eu/publications/emep-eea-guidebook-2023}.

\bibitem{emep2023si}
European Environment Agency, \emph{EMEP/EEA Guidebook 2023, gasoline
LDV PM emission factors by Euro class}, 2023.

\bibitem{emep2023nonexhaust}
European Environment Agency, \emph{EMEP/EEA Guidebook 2023, tyre,
brake and road wear (1.A.3.b.vi)}, update 2025.

\bibitem{beddows2021}
D.~Beddows and R.~M. Harrison, \emph{PM\textsubscript{10} and
PM\textsubscript{2.5} emission factors for non-exhaust particles
from road vehicles}, Atmospheric Environment, vol.~244, p.~117886,
2021.

\bibitem{hagino2016}
H.~Hagino, \emph{Laboratory testing of airborne brake-wear particle
emissions using a dynamometer system under urban city driving cycles
and vehicle braking maneuvers}, Atmospheric Environment, vol.~131,
pp.~269--278, 2016.

\end{thebibliography}

\end{document}
